% Generated by task tex-libraries from dossiers/lib-others.md, section C (and F.3 for the new
% pin). Snippets are copied from the dossier, which copied them from the sources.
\begingroup\sloppy
\section{Clean: circuits, tables, buses, and why its balance is unsatisfiable over \texorpdfstring{$\K$}{K}}\label{sec:gen-lib-clean}

Clean is a library for writing zkVM circuits (AIR tables and the buses between them) with proofs
that the constraints imply a specification. leanISA (the arithmetization) is written in it; the
proof system consumes its \emph{tables} through the spine's adaptor, not its soundness statement.
Every snippet of this section is from Clean at the old pin \code{93c9d1ef}, read with
\code{git show 93c9d1ef:<path>} under \file{.lake/}\allowbreak\file{packages/}\allowbreak\file{Clean}. At this pin Clean is
\textbf{not} a Lean \code{module} library, which is why
\file{LeanerVM/}\allowbreak\file{Parameters/}\allowbreak\file{CleanField.lean} is a plain file (\cref{sec:gen-lib-comppoly-leanervm}).
\Cref{sec:gen-lib-clean-new} says what the new pin \code{42fe4b26} changes.

\subsection{The circuit language}\label{sec:gen-lib-clean-circuit}

\begin{define}{\lean{FiniteField}, \lean{Expression}, \lean{Environment}, \lean{Operation}, \lean{Circuit}}
Every Clean circuit is generic over a field with Clean's \lean{FiniteField} interface: a Mathlib
field together with a reading \lean{val} into the natural numbers, its inverse \lean{fromNat},
and a \lean{size}. An \lean{Expression} is a polynomial expression over the variables of a row.
An \lean{Environment} assigns the variables (\lean{get}) and names the side tables a
specification may mention (\lean{data}, of type
\lean{ProverData F = String → (n : ℕ) → Array (Vector F n)}, \file{Expression.lean:20-21}). A
circuit emits \lean{Operation}s: a fresh witness variable, a polynomial asserted to be zero, a
lookup, a bus interaction, or a nested circuit; a \lean{Circuit} is a state monad producing
operations from a variable offset.
\end{define}

\begin{leancode}
class FiniteField (F : Type) extends Field F where
  /-- Canonical embedding of field elements into natural numbers. -/
  val : F → ℕ
  /-- Inverse of `val`: interpret a natural number (below the field size) as a field
  element. For prime fields this is `Nat.cast`; for binary fields it interprets the
  binary digits as polynomial coefficients. Note that `Nat.cast` would be wrong there:
  it reduces via the characteristic, collapsing `GF(2^n)` to `{0, 1}`. -/
  fromNat : ℕ → F
  ...
  size : ℕ
\end{leancode}

\src{\file{Clean/}\allowbreak\file{Utils/}\allowbreak\file{FiniteField.lean:23-45}, Clean \code{93c9d1ef}}


\begin{leancode}
inductive Expression (F : Type) where
  | var : Variable F -> Expression F
  | const : F -> Expression F
  | add : Expression F -> Expression F -> Expression F
  | mul : Expression F -> Expression F -> Expression F
\end{leancode}

\src{\file{Clean/}\allowbreak\file{Circuit/}\allowbreak\file{Expression.lean:12-16}, Clean \code{93c9d1ef}}


\begin{leancode}
structure Environment (F : Type) where
  /-- Assignment of a circuit's variables to field elements -/
  get : ℕ → F
  ...
  data : ProverData F
\end{leancode}

\src{\file{Clean/}\allowbreak\file{Circuit/}\allowbreak\file{Expression.lean:46-52}, Clean \code{93c9d1ef}}


\begin{leancode}
inductive Operation (F : Type) [FiniteField F] where
  | witness : (m : ℕ) → (compute : WitgenIR F m) → Operation F
  | assert : Expression F → Operation F
  | lookup : Lookup F → Operation F
  | interact : AbstractInteraction F → Operation F
  | subcircuit : {n : ℕ} → Subcircuit F n → Operation F
...
def Operations (F : Type) [FiniteField F] := List (Operation F)
\end{leancode}

\src{\file{Clean/}\allowbreak\file{Circuit/}\allowbreak\file{Operations.lean:315-320} and \file{:355}, Clean \code{93c9d1ef}}


\begin{leancode}
def Circuit (F : Type) [FiniteField F] (α : Type) := ℕ → α × List (Operation F)
\end{leancode}

\src{\file{Clean/}\allowbreak\file{Circuit/}\allowbreak\file{Basic.lean:29}, Clean \code{93c9d1ef}}

\subsection{Components, tables, ensembles}\label{sec:gen-lib-clean-tables}

\begin{define}{\lean{Component}, \lean{Table}, \lean{Ensemble}, \lean{EnsembleWitness}}
A \lean{Component} is one circuit checked on every row (\lean{GeneralFormalCircuit} bundles the
circuit with its assumptions, specification and soundness proof). A \lean{Table} is a component
with its concrete rows. An \lean{Ensemble} is a list of components and the channels (buses)
between them; its witness, an \lean{EnsembleWitness}, is one table per component, the shared
data and the public input.
\end{define}

\begin{leancode}
structure Component (F : Type) [FiniteField F] where
  {Input : TypeMap} {Output : TypeMap}
  [provableInput : ProvableType Input] [provableOutput : ProvableType Output]
  circuit : GeneralFormalCircuit F Input Output
\end{leancode}

\src{\file{Clean/}\allowbreak\file{Air/}\allowbreak\file{FlatComponent.lean:7-15}, Clean \code{93c9d1ef}}


\begin{leancode}
structure Table (F : Type) [FiniteField F] where
  component : Component F
  width : ℕ
  table : List (Array F)
  data : ProverData F
  uniform_width : ∀ row ∈ table, row.size = width
\end{leancode}

\src{\file{Clean/}\allowbreak\file{Air/}\allowbreak\file{FlatComponent.lean:150-156}, Clean \code{93c9d1ef}}


\begin{leancode}
structure Ensemble (F : Type) [FiniteField F] (PublicIO : TypeMap) [ProvableType PublicIO] where
  tables : List (Component F)
  channels : List (RawChannel F)
  -- TODO: the verifier shouldn't be treated as a "circuit", and possibly shouldn't even be on here
  verifier : GeneralFormalCircuit F PublicIO unit := .empty F PublicIO
  ...
structure EnsembleWitness (ens : Ensemble F PublicIO) where
  tables : List (Table F)
  data : ProverData F
  publicInput : PublicIO F
  same_length : ens.tables.length = tables.length
  same_circuits : ∀ i (hi : i < ens.tables.length), ens.tables[i] = tables[i].component
  same_data : ∀ table ∈ tables, table.data = data
\end{leancode}

\src{\file{Clean/}\allowbreak\file{Air/}\allowbreak\file{FlatEnsemble.lean:11-25}, Clean \code{93c9d1ef}}

\paragraph{Why \texorpdfstring{\texttt{EnsembleWitness}}{EnsembleWitness} lives in \texorpdfstring{\texttt{Type 1}}{Type 1}.}
A \lean{Table} stores its \lean{Component}, which stores its row types
\lean{Input Output : TypeMap}, and \lean{TypeMap := Type → Type}
(\file{Clean/}\allowbreak\file{Circuit/}\allowbreak\file{CircuitType.lean:10}) is itself a \lean{Type 1}. A structure with a field in
\lean{Type 1} is in \lean{Type 1} (probe \code{CleanBalance}, output lines 1--4:
\lean{Air.}\allowbreak\lean{Flat.}\allowbreak\lean{Component : (F : Type) → [FiniteField F] → Type 1}, likewise \lean{Table},
\lean{Ensemble}, \lean{EnsembleWitness}). ArkLib's statements and witnesses are in \lean{Type},
so an \lean{EnsembleWitness} cannot be the proof system's witness; the blueprint draws exactly
this consequence (\file{protocol-blueprint.md:367-369}, \file{:582-586}) and makes the stack
\lean{q : Column μ}, a \lean{Type}, the witness, with a pointwise transport along the adaptor.

\subsection{Channels and interactions}\label{sec:gen-lib-clean-channels}

\begin{define}{\lean{RawChannel}, \lean{Channel.toRaw}, \lean{Interaction}}
A raw channel is a name, a message arity, and two propositions indexed by the multiplicity: what
a receiver may assume (\lean{Guarantees}) and what a sender must prove (\lean{Requirements}). The
typed channel reads the \emph{direction from the sign} of the multiplicity: multiplicity $-1$ is a
pull, any other non-zero multiplicity a push. An interaction is a channel, a multiplicity and a
message: symbolic in a circuit (\lean{AbstractInteraction}), evaluated on a row
(\lean{Interaction}).
\end{define}

\begin{leancode}
structure RawChannel (F : Type) where
  name : String
  arity : ℕ
  Guarantees (mult : F) (message : Vector F arity) (data : ProverData F) : Prop
  Requirements (mult : F) (message : Vector F arity) (data : ProverData F) : Prop
\end{leancode}

\src{\file{Clean/}\allowbreak\file{Circuit/}\allowbreak\file{Channel.lean:15-25}, Clean \code{93c9d1ef}}


\begin{leancode}
def toRaw (channel : Channel F Message) : RawChannel F where
  name := channel.name
  arity := size Message
  Guarantees mult message data :=
    mult = -1 → channel.Guarantees (fromElements message) data
  Requirements mult message data :=
    mult ≠ -1 →
    mult ≠ 0 →
    channel.Guarantees (fromElements message) data
\end{leancode}

\src{\file{Channel.lean:35-45}, Clean \code{93c9d1ef}}


\begin{leancode}
structure AbstractInteraction (F : Type) where
  channel : RawChannel F
  mult : Expression F
  msg : Vector (Expression F) channel.arity
  assumeGuarantees : Bool
...
structure Interaction (F : Type) where
  channel : RawChannel F
  mult : F
  msg : Array F
  same_size : msg.size = channel.arity
  assumeGuarantees : Bool
\end{leancode}

\src{\file{Channel.lean:101-105} and \file{:305-310}, Clean \code{93c9d1ef}}

\subsection{Clean's ensemble statement and its soundness}\label{sec:gen-lib-clean-statement}

\begin{define}{\lean{balanceOf}, \lean{BalancedInteractions}, \lean{Ensemble.Statement}}
The balance of a message on a list of interactions is the sum, \emph{in the field}, of the
multiplicities of the interactions carrying it. A list is balanced when every message has balance
zero and a side condition bounds the number of interactions by the characteristic. The ensemble
statement for a public input says that some witness has that public input, satisfies every
table's constraints, and balances every channel; \lean{Soundness} says that the assumptions and
the statement imply the specification.
\end{define}

\begin{leancode}
/--
Balance of one element of an interaction list.
This is just multiplicity of the element when viewing the list as a multiset.
-/
def balanceOf (interactions : List (Interaction F)) (msg : Array F) : F :=
  interactions.filter (·.msg = msg) |>.map (·.mult) |>.sum

/--
Channel balance: for any message, the sum of multiplicities is 0.
We also require a side condition that ensures the interaction count does not overflow.
-/
def BalancedInteractions (interactions : List (Interaction F)) : Prop :=
  (interactions.length < ringChar F ∨ ringChar F = 0) ∧
  ∀ msg : Array F, balanceOf interactions msg = 0
\end{leancode}

\src{\file{Clean/}\allowbreak\file{Air/}\allowbreak\file{Balance.lean:13-26}, Clean \code{93c9d1ef}}


\begin{leancode}
abbrev BalancedChannel [DecidableEq F] {ens : Ensemble F PublicIO} (witness : EnsembleWitness ens)
    (channel : RawChannel F) : Prop :=
  BalancedInteractions (witness.allTablesWitness.interactionsWith channel)
...
def Statement (ens : Ensemble F PublicIO) (publicInput : PublicIO F) : Prop :=
  ∃ witness : EnsembleWitness ens,
    witness.publicInput = publicInput ∧
    witness.Constraints ∧
    witness.BalancedChannels

/-- Soundness: assumptions plus the raw statement imply the spec. -/
def Soundness (ens : Ensemble F PublicIO) (Assumptions Spec : PublicIO F → Prop) : Prop :=
  ∀ publicInput, Assumptions publicInput → ens.Statement publicInput → Spec publicInput
\end{leancode}

\src{\file{Clean/}\allowbreak\file{Air/}\allowbreak\file{FlatEnsemble.lean:334-337}, \file{:361-370}, Clean \code{93c9d1ef}}

The theorem that assembles per-table soundness into ensemble soundness has a proof that is plain
logic; the bus argument is inside \lean{TableSoundness} and the channel lemmas of
\file{Balance.lean} (for example \lean{exists_push_of_pull}, \file{:155-157}: every pull has a
matching push, derived from balance \emph{and} the side condition).

\begin{leancode}
theorem soundness_of_tableSoundness_and_specConsistency (ens : Ensemble F PublicIO)
  (Assumptions Spec : PublicIO F → Prop) :
  ens.TableSoundness →
  ens.AssumptionsConsistency Assumptions →
  ens.SpecConsistency Spec →
    ens.Soundness Assumptions Spec := by
\end{leancode}

\src{\file{Clean/}\allowbreak\file{Air/}\allowbreak\file{FlatEnsemble.lean:419-424}, Clean \code{93c9d1ef}}

\subsection{Why the balance argument is vacuous over a field of characteristic two}\label{sec:gen-lib-clean-vacuous}

The balance is a LogUp-style statement: multiplicities are \textbf{summed in the field}. That is
sound over a large prime field because a sum of at most $p - 1$ terms $\pm 1$ vanishes only if
the counts match; the side condition \lean{length < ringChar F} is that bound. Over $\K$,
\lean{ringChar K = 2}:
\begin{enumerate}
\item \textbf{Without the side condition the sum is unsound.} A message pushed twice
  (multiplicity $1$, twice) and never pulled has \lean{balanceOf} $= 1 + 1 = 0$. The probe
  \code{CleanBalance} proves both statements below (axioms: the standard three and
  \lean{[propext]} respectively):

\begin{leancode}
def pushedTwice : List (Interaction K) := [push 7, push 7]
theorem pushedTwice_balance (msg : Array K) : balanceOf pushedTwice msg = 0
theorem pushedTwice_not_perm :
    ¬ (pushedTwice.map (·.msg)).Perm ([] : List (Array K))
\end{leancode}

\src{probe \code{CleanBalance} (\file{probes/}\allowbreak\file{lib-others/}\allowbreak\file{CleanBalance.lean}), run at the old pins; \cref{ch:probes}}
  In addition $-1 = 1$ in $\K$ (\lean{example : (-1 : K) = 1 := by decide}), so
  \lean{Channel.toRaw} grants the pull guarantee to every multiplicity-$1$ push and makes every
  \lean{Requirements} vacuous (leanerVM issue \#16, table row 1, with kernel-checked examples in
  \file{tests/}\allowbreak\file{LeanerVMTests/}\allowbreak\file{Arithmetization/}\allowbreak\file{Channels.lean}).
\item \textbf{With the side condition, the relation holds of no real bus.} \lean{length < 2}
  admits at most one interaction per channel, and a balanced singleton has multiplicity $0$. The
  probe \code{CleanBalance} proves all of the following on the standard axioms:

\begin{leancode}
theorem pushedTwice_not_balanced : ¬ BalancedInteractions pushedTwice
theorem length_le_one_of_balanced {l : List (Interaction K)} (h : BalancedInteractions l) :
    l.length ≤ 1
theorem mult_eq_zero_of_balanced_singleton {i : Interaction K}
    (h : BalancedInteractions [i]) : i.mult = 0
theorem honest_pair_not_balanced : ¬ BalancedInteractions [push 7, pull 7]
theorem statement_forces_at_most_one_interaction {PublicIO : TypeMap} [ProvableType PublicIO]
    (ens : Ensemble K PublicIO) (w : EnsembleWitness ens) (h : w.BalancedChannels)
    (c : RawChannel K) (hc : c ∈ ens.channels) :
    (w.allTablesWitness.interactionsWith c).length ≤ 1
\end{leancode}

\src{probe \code{CleanBalance} (\file{probes/}\allowbreak\file{lib-others/}\allowbreak\file{CleanBalance.lean}), run at the old pins; \cref{ch:probes}}
  So over $\K$ \lean{Ensemble.Statement} is false for every ensemble whose honest witness puts two
  interactions on a channel, including the smallest honest bus (one push, one matching pull):
  Clean's \lean{Soundness} is then \emph{vacuously} true and its completeness false. That is the
  precise sense of \enquote{vacuous}: the countermodel of item 1 is excluded, but only by
  excluding every real witness too.
\end{enumerate}
The blueprint's relation-ladder row (\file{protocol-blueprint.md:359}) says \enquote{Over $\K$ the
side condition fails for every real ensemble and a tuple pushed twice and never pulled balances}.
The first half is right; the second is true of the field sum, not of
\lean{BalancedInteractions}, which rejects that list (\lean{pushedTwice_not_balanced}): the
finding \cref{f:lib-ladder-conflates}.

\subsection{State of the tracking items}\label{sec:gen-lib-clean-tracking}

Read with \code{gh} on 2026-09-30.

{\small
\begin{longtable}{@{}p{0.13\linewidth}p{0.22\linewidth}p{0.59\linewidth}@{}}
\caption{The issues and pull requests on Clean's bus balance and fixed columns.}\label{tab:gen-lib-clean-tracking}\\
\toprule
Item & State & What it is \\
\midrule
\endfirsthead
\toprule
Item & State & What it is \\
\midrule
\endhead
\bottomrule
\endfoot
leanerVM \#16 & open (updated 2026-09-15) & \enquote{leanISA: Clean's bus argument is vacuous over K,
  PR Clean for direction tags and multiset balance}: the request; tabulates \lean{Channel.toRaw},
  \lean{emit}, \lean{BalancedInteractions}, and the \lean{[Fact (ringChar F ≠ 2)]} lemmas as
  degenerate over $\K$. \\
leanerVM \#20 & open (updated 2026-09-24) & the Clean-side roadmap: opt-in \lean{DirectedChannel}
  (direction as a tag in the message), \lean{BalanceModel.multiset} (permutation, counted in
  $\N$), \lean{Ensemble.}\allowbreak\lean{StatementWith model}; Layers 0--4 on Clean \#464, Layer 5 (leanerVM
  adoption) open. \\
leanerVM \#23 & open (updated 2026-09-17) & fold the three named fixed-column hypotheses of
  \lean{SatisfiedBy} into Clean fixed columns and proof-committed data once Clean \#446 lands. \\
Clean \#446 & open, \textbf{draft} (updated 2026-08-16) & \enquote{Add fixed columns and sound
  prover data to Flat AIR}. \\
Clean \#452 & open \textbf{issue} (not a pull request) & \enquote{Model verifier-checked
  interaction capacity in Flat AIR soundness}. \\
Clean \#464 & open, ready for review (updated 2026-09-24) & \enquote{Bus balance over binary
  fields: opt-in directed channels and explicit balance models}; legacy statements unchanged;
  includes \lean{legacy_rejects_every_run} over \lean{F 2}. \\
Clean \#466 & open (updated 2026-09-24) & \enquote{interpret expressions as bounded-degree
  polynomials} (the spine's hole \enquote{Clean components as polynomials},
  \file{protocol-blueprint.md:604}). \\
Clean \#474 & open (updated 2026-09-29) & \enquote{preserve balance proof across Lean 4.34
  normalization}: the \textbf{new pin \code{42fe4b26}} is this branch merged with Clean
  \code{main}; it ports proofs only. At \code{42fe4b26} \lean{balanceOf} and
  \lean{BalancedInteractions} are \textbf{unchanged} (\file{Clean/}\allowbreak\file{Air/}\allowbreak\file{Balance.lean:21-29},
  identical text), so the vacuity stands at the new pin. \\
\end{longtable}}

\subsection{What leanISA does instead, and when the workaround retires}\label{sec:gen-lib-clean-workaround}

leanISA never consumes \lean{Ensemble.Statement}. Its relation \lean{SatisfiedBy} states balance
as a permutation of message lists, counted in $\N$, with direction given by \emph{which channel
of a pair} an interaction is on:

\begin{leancode}
/-- Pushed and pulled messages of a channel pair form the same multiset (specification §5.1):
a permutation of message lists, counted in `ℕ`, never a field sum (acceptance test 14). -/
def BalancedPair (w : EnsembleWitness (leanIsaEnsemble prog)) (pull push : RawChannel K) :
    Prop :=
  (messagesOn w push).Perm (messagesOn w pull)
\end{leancode}

\src{\file{LeanerVM/}\allowbreak\file{Arithmetization/}\allowbreak\file{Statement.lean:227-231}, leanerVM \code{b435631}}

\lean{BalancedPair} is used three times in \lean{SatisfiedBy} (\file{Statement.lean:320-324}:
\lean{state_balanced}, \lean{mem_balanced}, \lean{bytecode_balanced}). Every interaction has
multiplicity $1$, so the message multiset is the interaction multiset (\file{Statement.lean:57-64};
\file{Channels.lean:58-72} for the channel pairs). \lean{List.Perm} (core: \lean{l₁ ~ l₂}, one list
is a reordering of the other) is exact multiset equality with no characteristic condition.

\paragraph{Retirement condition.} \lean{BalancedPair} and the channel-pair encoding of direction
(\lean{Direction}, \lean{channelDir}) are deleted when (1) Clean \#464 is merged into Clean
\code{main}, (2) leanerVM's pin reaches a Clean revision containing it, and (3) leanerVM's
channels are rebuilt as \lean{DirectedChannel}s and \lean{SatisfiedBy} is stated through
\lean{Ensemble.}\allowbreak\lean{StatementWith BalanceModel.}\allowbreak\lean{multiset} (leanerVM \#20, Layer 5, \enquote{a separate
pull request}; \file{Channels.lean:71-72}). Until then the workaround is \textbf{a deviation
forced by an upstream library}, and its correctness is by inspection of a four-line definition.
Neither \#446 nor \#452 is required for the balance repair (issue \#20: \enquote{\#446
(independent)}).

\subsection{At the new pin \texorpdfstring{\texttt{42fe4b26}}{42fe4b26}}\label{sec:gen-lib-clean-new}

Clean is a \code{module} library (\#460), so \file{LeanerVM/}\allowbreak\file{Parameters/}\allowbreak\file{CleanField.lean} became a
\code{module} at \code{144c5aa} and the \enquote{Clean bridge is plain} convention (leanISA status
finding C8) retires; the balance proofs were ported (\#474); \lean{balanceOf} and
\lean{BalancedInteractions} are unchanged, and \#464 (directed channels) is not in the pin (a
\code{git grep DirectedChannel} of \code{42fe4b26} finds nothing), so the vacuity over $\K$ and
leanISA's workaround stand as described above.

\par\endgroup